% Options for packages loaded elsewhere
\PassOptionsToPackage{unicode}{hyperref}
\PassOptionsToPackage{hyphens}{url}
%
\documentclass[
]{article}
\usepackage{amsmath,amssymb}
\usepackage{iftex}
\ifPDFTeX
  \usepackage[T1]{fontenc}
  \usepackage[utf8]{inputenc}
  \usepackage{textcomp} % provide euro and other symbols
\else % if luatex or xetex
  \usepackage{unicode-math} % this also loads fontspec
  \defaultfontfeatures{Scale=MatchLowercase}
  \defaultfontfeatures[\rmfamily]{Ligatures=TeX,Scale=1}
\fi
\usepackage{lmodern}
\ifPDFTeX\else
  % xetex/luatex font selection
\fi
% Use upquote if available, for straight quotes in verbatim environments
\IfFileExists{upquote.sty}{\usepackage{upquote}}{}
\IfFileExists{microtype.sty}{% use microtype if available
  \usepackage[]{microtype}
  \UseMicrotypeSet[protrusion]{basicmath} % disable protrusion for tt fonts
}{}
\makeatletter
\@ifundefined{KOMAClassName}{% if non-KOMA class
  \IfFileExists{parskip.sty}{%
    \usepackage{parskip}
  }{% else
    \setlength{\parindent}{0pt}
    \setlength{\parskip}{6pt plus 2pt minus 1pt}}
}{% if KOMA class
  \KOMAoptions{parskip=half}}
\makeatother
\usepackage{xcolor}
\usepackage[margin=1in]{geometry}
\usepackage{graphicx}
\makeatletter
\def\maxwidth{\ifdim\Gin@nat@width>\linewidth\linewidth\else\Gin@nat@width\fi}
\def\maxheight{\ifdim\Gin@nat@height>\textheight\textheight\else\Gin@nat@height\fi}
\makeatother
% Scale images if necessary, so that they will not overflow the page
% margins by default, and it is still possible to overwrite the defaults
% using explicit options in \includegraphics[width, height, ...]{}
\setkeys{Gin}{width=\maxwidth,height=\maxheight,keepaspectratio}
% Set default figure placement to htbp
\makeatletter
\def\fps@figure{htbp}
\makeatother
\setlength{\emergencystretch}{3em} % prevent overfull lines
\providecommand{\tightlist}{%
  \setlength{\itemsep}{0pt}\setlength{\parskip}{0pt}}
\setcounter{secnumdepth}{5}
\renewcommand{\contentsname}{Contenido}

\usepackage{float}
\usepackage{caption}
\usepackage{xcolor}

\captionsetup{
  font={bf,small},
  labelfont={bf},
  justification=centering,
  position=top,
  singlelinecheck=false
}

\captionsetup[table]{
  position=top,
  skip=5pt
}

\captionsetup[figure]{
  position=above,
  skip=5pt
}
\ifLuaTeX
  \usepackage{selnolig}  % disable illegal ligatures
\fi
\usepackage{bookmark}
\IfFileExists{xurl.sty}{\usepackage{xurl}}{} % add URL line breaks if available
\urlstyle{same}
\hypersetup{
  hidelinks,
  pdfcreator={LaTeX via pandoc}}

\author{}
\date{\vspace{-2.5em}}

\begin{document}

\begin{titlepage}
\centering

\vspace*{1.5cm}

% --- Título del proyecto ---
{\Huge \textbf{Métrica RAPM: Evaluación del impacto ofensivo y defensivo de los jugadores de la Liga Nacional de Básquetbol en Argentina}\par}

\vspace{1.5cm}

% --- Subtítulo ---
{\LARGE \textbf{Tesina de grado}\par}

\vspace{2cm}

% --- Datos del autor y directores ---
{\Large Estudiante: Simón Pedro Gazze\par}
{\Large Director/a: Andrés Sosa\par}
{\Large Co-director/a: Cristina Cuesta\par}

\vspace{2cm}

% --- Fecha ---
{\Large Fecha: 2 de noviembre de 2025\par}

\vfill

% --- Logo centrado ---
\includegraphics[height=2.8cm]{Logo-unr.jpg}

\end{titlepage}

\newpage

\renewcommand{\contentsname}{Contenido}
\tableofcontents
\newpage

\section{Resumen}\label{resumen}

El crecimiento de la analítica deportiva ha impulsado el desarrollo de
métricas destinadas a evaluar con mayor precisión la contribución de los
jugadores al rendimiento de sus equipos. Particularmente en el
básquetbol, las estadísticas individuales derivadas del box score
presentan limitaciones para capturar el impacto real de los jugadores
sobre el resultado del juego. Una de las medidas desarrolladas con este
propósito es el \emph{plus-minus}, que registra la diferencia de puntos
generada mientras un jugador permanece en cancha. Sin embargo, esta
métrica presenta dificultades para aislar el efecto propio de cada
jugador respecto del contexto en el que participa. En respuesta a esta
problemática surge el \emph{Regularized Adjusted Plus Minus}
(\emph{RAPM}), una medida basada en modelos estadísticos que estima el
impacto ofensivo y defensivo de cada jugador controlando simultáneamente
la presencia de compañeros y rivales en cancha.

El objetivo de este trabajo es estimar el impacto individual de los
jugadores de la Liga Nacional de Básquetbol de Argentina durante la
temporada 2024/2025 mediante modelos RAPM, construyendo rankings que
permitan identificar y comparar sus contribuciones ofensivas y
defensivas. Asimismo, se analizan distintas alternativas metodológicas
para abordar problemas recurrentes en este tipo de modelos,
particularmente la multicolinealidad entre jugadores y la presencia de
coeficientes extremos asociados a jugadores con escasa participación en
cancha.

Para esto, se construyó una base de datos \emph{play-by-play} a partir
de información obtenida mediante técnicas de \emph{web scraping} sobre
los partidos de la temporada regular. Posteriormente, las acciones
registradas fueron agrupadas en posesiones, que constituyen la unidad de
análisis utilizada para la estimación de las métricas propuestas.
Inicialmente, se ajustaron modelos de regresión lineal múltiple,
incorporando distintas técnicas de regularización. Posteriormente, se
implementaron modelos de regresión logística multinomial que representan
de manera más adecuada la naturaleza discreta de la variable respuesta.
A partir de las probabilidades estimadas por estos modelos se
construyeron métricas basadas en puntos esperados por posesión
(\emph{EPTS}). Finalmente, se propuso una ponderación a dicha métrica
denominada \emph{wEPTS}, que incorpora la participación de cada jugador
relativa a las posesiones de su equipo con el objetivo de penalizar a
los jugadores con poco tiempo de juego. Los resultados obtenidos
muestran que esta métrica presentó un desempeño superior al de las
alternativas propuestas, de acuerdo con los distintos criterios de
validación utilizados.

Este estudio constituye la primera aplicación de métricas \emph{RAPM} en
la Liga Nacional de Básquetbol de Argentina y aporta una metodología
reproducible para la evaluación objetiva del rendimiento individual de
los jugadores, contribuyendo al desarrollo de herramientas de analítica
deportiva en el contexto del básquetbol argentino.

\newpage

\section{Introducción}\label{introducciuxf3n}

\subsection{Sports Analytics}\label{sports-analytics}

\emph{Sports Analytics} o analítica en el deporte es un término que
refiere a la aplicación de métodos estadísticos para describir, explicar
y/o predecir fenómenos vinculados al rendimiento deportivo, tanto a
nivel individual como colectivo con el objetivo de facilitar la toma de
decisiones de entrenadores o \emph{staff} técnico antes, durante y
después de los partidos.

A pesar de que las primeras aplicaciones al análisis de datos al deporte
se remontan a mediados del siglo XX, el impulso más significativo
ocurrió a partir del año 2000, en particular en el béisbol, con la
popularización del enfoque conocido como \emph{sabermetrics}\footnote{Término
  acuñado por Bill James (1980) para referirse al análisis estadístico
  del béisbol, basado en la recopilación y estudio sistemático de datos
  con el fin de comprender y cuantificar el desempeño de los jugadores y
  equipos. El nombre proviene de SABR (\emph{Society for American
  Baseball Research}).}. Desde entonces, la práctica se extendió a otros
deportes como el básquetbol, el fútbol o el tenis; dando lugar a una
nueva cultura de toma de decisiones basada en evidencia. Hoy en día, los
equipos profesionales, las ligas y las federaciones emplean
departamentos especializados en análisis de datos para optimizar
estrategias de juego, evaluar el rendimiento de jugadores, prevenir
lesiones y gestionar recursos económicos de manera más eficiente.

El \textbf{básquetbol} es uno de los deportes que más ha incorporado el
análisis estadístico en su evolución reciente. Desde las primeras
métricas derivadas del \emph{box score} (resumen estructurado de los
resultados de un partido), el desarrollo de bases de datos
\emph{play-by-play} (registro de cada una de las acciones del partido) a
partir de la década del 2000 hasta la incorporación más reciente de
tecnologías de \emph{player tracking}\footnote{El \emph{player tracking}
  es un sistema tecnológico que utiliza cámaras y software de visión
  artificial en los estadios para registrar los movimientos de jugadores
  y pelotas en tiempo real en la cancha, produciendo datos detallados
  sobre su rendimiento.}. A partir de esta información, se han
desarrollado distintas herramientas estadísticas para poder estudiar la
eficiencia de los jugadores, la eficiencia de los tiros al aro, el
desempeño de los equipos, la importancia en la creación de espacios,
entre otras cosas.

\subsection{Medidas de desempeño para jugadores de
básquetbol}\label{medidas-de-desempeuxf1o-para-jugadores-de-buxe1squetbol}

Desde la consolidación del \emph{box-score} en el básquetbol
profesional, durante la década de 1950, empezaron a calcularse las
primeras medidas de desempeño para los jugadores, como pueden ser los
puntos, rebotes, asistencias, etc. Durante mucho tiempo estas métricas
constituyeron una de las formas más tradicionales de evaluar el
rendimiento individual en el básquetbol. Su amplia disponibilidad y
fácil interpretación las convirtieron en el principal insumo para
comparar jugadores, otorgar premios o negociar contratos. Estas medidas
suelen agruparse en un conjunto de estadísticas denominadas
\emph{bottom-up}, las cuales se calculan en función de las acciones
individuales de los jugadores, es decir, si un jugador realiza una
acción que se asocia con un resultado positivo para el equipo, la
calificación de ese jugador aumenta, mientras que la calificación de los
demás jugadores, que no participaron en la jugada, permanece sin
cambios. Pero estas medidas presentan algunas limitaciones importantes,
ya que, sólo capturan una parte de los eventos relevantes del partido
(acciones como una defensa asfixiante, buenas rotaciones defensivas o
pases rápidos para iniciar una posesión no están contempladas) y no
necesariamente reflejan el impacto real de un jugador en las victorias
de su equipo, que es lo que importa a fin de cuentas.

Por otra parte, las medidas \emph{top-down} se basan en el rendimiento
del equipo en su conjunto, y el crédito por dicho rendimiento se
distribuye entre los jugadores que estuvieron involucrados en el
partido, sin importar qué acciones puntuales hayan realizado.

Con la idea de que lo más importante en un partido es \textbf{ganar}, y
entendiendo que las estadísticas tradicionales son solo una medida
imperfecta de muchas de las contribuciones que realizan los jugadores a
la victoria, en este trabajo buscaremos ``pensar más allá del \emph{box
score}'' y nos centraremos en estadísticas \emph{top-down}, más
precisamente en la medida \emph{Plus/Minus}. Esta estadística registra
los cambios que se producen en el marcador durante los minutos en que el
jugador de interés está en cancha, partiendo de la idea de que los
equipos deberían rendir mejor (reflejado en un \emph{Plus/Minus} más
alto) cuando sus buenos jugadores están en cancha que cuando no lo
están. Sin embargo esta métrica posee una deficiencia clave: \textbf{la
calificación de cada jugador va a depender en gran medida de la calidad
de sus compañeros en el campo}. Por ejemplo: un jugador de rol dentro de
un gran equipo va a tener un mayor \emph{Plus/Minus} que una
superestrella en un equipo con mal desempeño.

\subsection{RAPM (Regularized Adjusted Plus
Minus)}\label{rapm-regularized-adjusted-plus-minus}

La deficiencia anteriormente mencionada se ha intentado solucionar
implementando \textbf{modelos matemáticos} que incorporen los efectos
tanto de los compañeros presentes en cancha como de los rivales en ese
momento, buscando aislar el efecto real de los jugadores en la cancha.
El primero en realizar una publicación sobre el tema fue Dan Rosenbaum
(2004), el cual planteó un \textbf{modelo de regresión lineal múltiple}
en base a observaciones provenientes de partidos de las temporadas
2002-2003 y 2003-2004 de la National Basketball Association o
\emph{NBA}. Como unidad de análisis se utilizaron segmentos de partidos
donde no ocurrian sustituciones, la variable respuesta fue la diferencia
en el marcador entre el equipo local y el visitante en ese segmento; y
como variables explicativas se postularon variables dicotómicas para
cada uno de los jugadores que hayan jugado al menos un partido en esas
temporadas. Finalmente las métricas de desempeño se correspondían con
los coeficientes estimados que acompañaban a las ``Dummies'' de cada
jugador, dichas estimaciones se realizaron mediante la técnica de
Mínimos Cuadrados Ordinarios. Estas estadísticas se conocen como
\emph{Adjusted Plus Minus} (APM).

A pesar de que estas métricas presentaban una solución innovadora a la
problematica inicial, Rosenbaum señaló que las estimaciones obtenidas no
resultaron muy precisas. El modelo propuesto presentaba estimaciones con
mucho error, debido a un claro problema de \textbf{multicolinealidad} ya
que muchos de los jugadores compartían una gran cantidad de minutos
juntos en cancha, y se necesitaba de una gran cantidad observaciones de
varias temporadas para que el modelo pudiera separar el efecto propio de
cada jugador.

Para solucionar está problemática, se plantea la métrica RAPM
introducida por Sill (2010) en la cual se aborda el problema de la
Multicolinealidad mediante estimaciones de los parametros realizadas con
Métodos de Regularización, particularmente utilizando la
\textbf{Regresión Ridge}. Esté método de regularización es el más
abordado a lo largo de toda la literatura relevante sobre RAPM hasta la
fecha.

Otro de los problemas recurrentes en estos modelos esta relacionado con
la inclusión de jugadores que participan muy pocos minutos en cancha
(\emph{LTPs Players}). Cuando un jugador juega muy poco, el modelo solo
tiene unas pocas observaciones (segmentos) para evaluarlo, y si
particularmente en esas observaciones el equipo tiene un muy buen
desempeño (probablemente por casualidad); entonces el modelo no puede
diferenciar si el efecto es propio del jugador y estima coeficientes
extremos. Para solucionar estos casos se han contemplado las siguientes
alternativas: 1)excluir a los jugadores con pocos minutos dentro de las
temporadas (por ejemplo: Rosembaum (2004) excluye a lo jugadores con
menos de 250 minutos, Ilardi (2007) excluyen a lo jugadores con menos de
400 minutos, etc), 2) utilizar una gran cantidad de observaciones
(Ilardi \& Barzilai (2008) incorpora información de 5 temporadas), 3)
plantear métodos de estimación de los parámetros que penalicen a este
tipo de jugadores.

En los últimos años, muchos investigadores han publicado trabajos
orientados a mejorar los resultados de las métricas RAPM.
Particularmente en este trabajo usaremos como referencia central al
artículo titulado: \textbf{Lasso multinomial performance indicators for
in‑play basketball data (Damoulaki et.al , 2025)}, el cual presenta
diferencias importantes respecto de los estudios previos. Dicho artículo
utiliza \textbf{posesiones}\footnote{Una posesión comienza (Kubatko et
  al., 2007) cuando un equipo obtiene el control de la pelota y termina
  cuando ese equipo cede el control de la misma al equipo rival.} como
unidad de análisis, \textbf{puntos} en esas posesiones como variable
respuesta y emplea \textbf{modelos logísticos multinomiales} que
representan de manera más adecuada la naturaleza de la respuesta.

\section{Motivación}\label{motivaciuxf3n}

El creciente interés por la analítica deportiva que se ha evidenciado en
los últimos años en distintos paises del mundo, ha empezado a captar la
atención de las entidades deportivas en Argentina. Particularmente este
año la Confederación Argentina de Básquetbol (CABB) y la Asociación de
Clubes (AdC) anunciaron en octubre la ampliación de su alianza con
Catapult (empresa australiana de tecnología deportiva), integrando
equipamiento de \emph{wearables} (GPS) y software de videoanálisis para
las 19 franquicias de la Liga Nacional, las 34 de la Liga Argentina y
las selecciones nacionales en todas las ramas. Mientras que en el ámbito
educativo, el Instituto CAB (creado en 2023) realizó un convenio con
Sports Data Campus (España) para ofrecer, en agosto de este año, el
primer curso de \emph{Big Data aplicado al básquetbol} en Argentina, con
el objetivo de que entrenadores y \emph{staff} técnico colaboren en
proyectos reales de la CABB, aportando soluciones basadas en datos sobre
rendimiento, \emph{scouting} y optimización de procesos\footnote{Confederación
  Argentina de Básquetbol (2025). \emph{La Confederación Argentina de
  Básquet y Sports Data Campus firman un acuerdo de colaboración}.
  \href{https://www.argentina.basketball/ver/noticia/la-confederacion-argentina-de-basquet-y-sports-data-campus-firman-un-acuerdo-de-colaboracion}{Link
  al artículo}.}.

Estas cuestiones previamente mencionadas dan indicio de que estamos ante
un momento ideal para profundizar en el estudio y la aplicación del
análisis de datos en el básquetbol. Se trata de un área que comienza a
reconocerse como un componente fundamental tanto para mejorar el
rendimiento deportivo de los equipos como para impulsar el desarrollo
del negocio en torno al deporte. Sin embargo, el potencial de la
analítica aplicada al básquetbol en Argentina aún pareciera encontrarse
lejos de estar plenamente aprovechado.

En este marco, y con el fin de caracterizar el estado de situación del
análisis de datos en el básquet argentino, se relevaron los proyectos e
iniciativas actualmente activos en el país vinculados a la estadística y
la tecnología aplicada al deporte.

\begin{itemize}
\item
  \emph{CHAS All Stats}, dirigido por Facundo Salas y Sebastián Fiol,
  ambos entrenadores de básquetbol Nivel 3 Eneba, especialistas en
  estadísticas avanzadas y analistas de datos en equipos profesionales.
  Ofrecen servicios de análisis táctico y brindan cursos de estadísticas
  avanzadas.
\item
  \emph{Básquet Advance}, dirigido por el entrenador Cristian Sánchez,
  ofrece informes pre/post partido con estadísticas avanzadas y
  tendencias clave para entrenadores y scouts.
\end{itemize}

Ambos proyectos ofrecen herramientas muy importantes y de gran valor
tanto para los equipos como para el desarrollo del análisis de datos en
el deporte. Pero particularmente en ambos casos, los que llevan adelante
estos proyectos son entrenadores de básquetbol, los cuales a pesar de
que pueden contar con una sólida formación técnica, no son profesionales
de ciencias estadística o matemática.

Por este motivo, resulta fundamental la incorporación de profesionales
en carreras de ciencias exactas, que trabajen en conjunto con los
entrenadores y cuerpos técnicos, para de esta manera poder ofrecer
nuevos análisis más rigurosos y profundos, así como mejorar las
herramientas actualmente disponibles. Esta correspondencia entre el
conocimiento técnico-deportivo y el enfoque estadístico permitiría
potenciar el desarrollo del básquet argentino y contribuir a elevar su
nivel competitivo en un futuro.

\newpage

\section{Objetivos}\label{objetivos}

\subsection{Objetivo general}\label{objetivo-general}

El objetivo general de esta tesina es estimar el aporte individual de
los jugadores de la Liga Nacional de Básquetbol de Argentina durante la
temporada 2024/2025 mediante modelos RAPM, con el fin de construir un
ranking que refleje su impacto en el desempeño de sus equipos.

\subsection{Objetivos específicos}\label{objetivos-especuxedficos}

\begin{enumerate}
\def\labelenumi{\arabic{enumi})}
\item
  Generar una base de datos del tipo \emph{play-by-play} para la
  temporada 2024/2025 de la Liga Nacional de Básquetbol de Argentina,
  mediante técnicas de \emph{web scraping}.
\item
  Diferenciar el aporte ofensivo y defensivo de cada uno de los
  jugadores de la liga, obteniendo finalmente rankings que identifiquen
  a los mejores jugadores en cada rubro.
\item
  Desarrollar una metodología adecuada que permita solucionar el
  problema de los coeficientes extremos asociados a los jugadores con
  pocos minutos de juego (\emph{LTPs Players}) durante la temporada.
\item
  Comparar modelos en base a distintos criterios de evaluación, para
  poder determinar cual de ellos proporciona mejores métricas, debiendo
  resumir de manera consistente y precisa el aporte ofensivo y defensivo
  de cada jugador en relación al desempeño de su equipo.
\end{enumerate}

El objetivo general de esta tesina es estimar de manera adecuada el
aporte individual de los jugadores de la Liga Nacional de Básquetbol de
Argentina durante la temporada 2024/2025, comparando enfoques
tradicionales basados en modelos de regresión lineal múltiple con
modelos logísticos multinomiales que se ajusten mejor a la naturaleza
discreta de la variable respuesta, evaluando en cada caso el uso de
distintas técnicas de regularización para la estimación de los
coeficientes.

Comparar los resultados provistos por los distintos modelos en base a
criterios externos de evaluación, para poder determinar cual de ellos
proporciona mejores métricas.

El objetivo general de esta tesina es evaluar y comparar diferentes
enfoques de estimación de modelos RAPM para cuantificar el impacto
individual de los jugadores de la Liga Nacional de Básquetbol de
Argentina durante la temporada 2024/2025, mediante el contraste de
modelos lineales y logísticos multinomiales, y el uso de distintas
técnicas de regularización para la estimación de los coeficientes.

Generar rankings definitivos de la Liga Nacional como el resultado final
obtenido mediante el modelo óptimo, identificando a los mejores
jugadores de cada rubro para la temporada en estudio.

\section{Materiales}\label{materiales}

Una parte central de esta tesina estará dedicada a la construcción de
una base de datos \emph{play-by-play} correspondiente a la temporada
2024/2025 de la Liga Nacional de Básquetbol de Argentina (LNB). Este
tipo de base registra de manera secuencial todas las acciones ocurridas
dentro de un partido. Dado que no existe una fuente oficial que provea
los datos en formato estructurado y descargable para su análisis
estadístico, fue necesario diseñar y ejecutar un proceso de \emph{web
scraping} que permita recolectar, integrar y depurar la información de
manera sistemática y reproducible. Los datos necesarios se extrajeron
del \href{https://www.laliganacional.com.ar/laliga/}{Sitio Oficial de la
Liga Nacional de Básquet} , utilizando librerías de Python como
\texttt{Selenium} y \texttt{BeautifulSoup}.

Posteriormente, se realizó un preprocesamiento orientado a transformar
los eventos crudos en un \emph{dataset} apto para la aplicación de los
modelos mencionados anteriormente. Como resultado, se obtuvo una tabla
ordenada cronológicamente en la cual cada fila corresponde a una
posesión del partido, contabilizando los puntos anotados durante esa
posesión y los jugadores que formaron parte de la misma, diferenciando
sus roles ofensivos y defensivos. Debido a que bajo está metodología se
registran cada una de las posesiones del partido de manera secuencial,
se denominó a esta estructura final como: base de datos
\emph{poss-by-poss}.

\subsubsection{Generación de la base de datos a partir de técnicas de
Web
Scraping}\label{generaciuxf3n-de-la-base-de-datos-a-partir-de-tuxe9cnicas-de-web-scraping}

\paragraph{Listado de links de los
partidos}\label{listado-de-links-de-los-partidos}

El primer paso consistió en la identificación y recopilación de los
enlaces a las páginas webs que contienen las estadísticas individuales
de todos los partidos incluidos en el período de estudio (378 partidos
de temporada regular). Para ello, se accedió a la página de la Liga
Nacional,
\href{https://www.laliganacional.com.ar/laliga/fixture}{sección
Fixture}, la cual presenta una tabla dinámica con los encuentros
disputados y permite acceder a la información detallada de cada partido.

\begin{figure}

{\centering \includegraphics[width=0.7\linewidth]{img/imagen links partidos} 

}

\caption{Fixture dentro de la página web de la LNB.}\label{fig:unnamed-chunk-1}
\end{figure}

En la Figura 1 se muestra la interfaz del sitio web, donde se visualiza
el listado completo de partidos de la temporada para el rango de fechas
seleccionadas. Dado que no pudo automatizarse la obtención de cada uno
de estos links, se procedió a copiar manualmente el \emph{DOM}
(\emph{Document Object Model}) de la página y a pegarlo en un archivo de
texto.

--figura2--

A partir de dicho archivo de texto, se utiliza la biblioteca
BeautifulSoup de Python para recorrer la estructura del DOM,
identificando previamente los elementos que contienen los enlaces de
cada partido. Se observa en la Figura 2 que la información de cada
partido se encuentra dentro de las clases ``fila-tabla-calendarios
datos-estadisticos'', por lo cual se recorren estas clases obteniendo
los respectivos enlaces a las páginas de estadísticas de cada partido,
los nombres de los equipos participantes y el resultado del encuentro,
para cada uno de los 378 partidos disputados en la temporada regular.
Obteniendo finalmente un diccionario que contenga toda esta información.

Dicho diccionario constituyó el insumo básico para el proceso posterior
de scraping masivo, sobre los cuales se iterará para obtener así la
información correspondiente a cada una de las acciones registradas en
cada partido. Esta estrategia permitió, además, asegurar la
reproducibilidad del proceso, ya que el listado de partidos es eliminado
de la página web una vez terminada la temporada, quedando fijado y
documentado de manera explícita para su uso, en caso de ser requerido.

\paragraph{\texorpdfstring{Base de datos
\emph{play-by-play}}{Base de datos play-by-play}}\label{base-de-datos-play-by-play}

Una vez obtenido el listado completo de enlaces correspondientes a cada
uno de los partidos en el período de estudio, se procedió a la
extracción del registro de acciones del partido (\emph{play-by-play}).
Para esto, se accedió individualmente a la página de estadísticas de
cada encuentro, particularmente a la pestaña denominada ``en vivo'',
donde se presenta el detalle secuencial de los eventos.

Esta sección del sitio web contiene una tabla que se genera
dinámicamente y que registra, en orden cronológico, todas las acciones
ocurridas durante el partido.

\begin{figure}

{\centering \includegraphics[width=0.7\linewidth]{img/imagen play by play_1} 

}

\caption{Pestaña en vivo dentro de la página web de estadísticas de un partido}\label{fig:unnamed-chunk-2}
\end{figure}

En dicha tabla, cada acción del partido se encuentra representada como
un elemento que incluye información textual sobre el tipo de acción (por
ejemplo, lanzamiento convertido, falta, pérdida, etc), el jugador que
realiza dicha acción, el cuarto de juego, el tiempo restante en el
cuarto y , en caso de ser una acción donde se generan puntos, el
marcador parcial en ese momento del partido. Esta información fue
sistemáticamente extraída y almacenada en una estructura tabular,
conservando el orden temporal de las acciones.

Este procedimiento se aplicó de forma idéntica a todos los partidos de
la temporada analizada, utilizando como insumo los enlaces previamente
identificados. Finalmente, los registros de acciones obtenidos para cada
encuentro fueron concatenados en una única base de datos de
\emph{play-by-play}, incorporando un identificador de partido que
permite distinguir y vincular las acciones correspondientes a cada
juego.

Cabe destacar que el \emph{play-by-play} resultante constituye una
secuencia detallada de eventos, pero en su forma original no incluye una
identificación explícita del equipo asociado a cada acción, lo cual
resulta indispensable a la hora de poder determinar a todos los
jugadores presentes, tanto de manera ofensiva como defensiva, en cada
una de las acciones del partido.

\paragraph{Tabla de jugadores-equipos}\label{tabla-de-jugadores-equipos}

Para resolver este problema, se obtuvo dicha información a partir de los
datos provenientes del \emph{Box Score}, localizados en la pestaña
``Estadísticas'' de la página web del partido. Aquí encontraremos
información referida a cada uno de los jugadores anotados en la planilla
de ambos equipos.

\begin{figure}

{\centering \includegraphics[width=0.7\linewidth]{img/imagen box score} 

}

\caption{Ejemplo del Box Score de un partido}\label{fig:unnamed-chunk-3}
\end{figure}

En esta pestaña, la información se organiza en dos tablas con
estadísticas individuales de los jugadores, una correspondiente al
equipo local y otra al equipo visitante, las cuales aparecen en el mismo
orden en que se muestran los nombres de los equipos en la interfaz del
sitio web. A partir de esta estructura, se extrajo para cada jugador su
nombre completo ,el equipo al que pertenece y la cantidad de minutos
disputados en dicho partido, generándose así una tabla de
correspondencia jugador--equipo.

Esta tabla constituye un diccionario que permite asociar a cada jugador
con su respectivo equipo, y es utilizada posteriormente como insumo para
realizar dicha correspondencia en cada acción de la base
\emph{play-by-play}. A su vez, mediante la adición de la cantidad de
minutos en el partido se podrá realizar un filtrado de jugadores que
hayan participado menos de cierto umbral requerido de tiempo durante la
temporada, cuestión fundamental a tratar posteriormente en este trabajo.

\subsubsection{Preprocesamiento de los
datos}\label{preprocesamiento-de-los-datos}

\paragraph{Integración del box score con el
play-by-play}\label{integraciuxf3n-del-box-score-con-el-play-by-play}

Una vez construida la tabla jugadores-equipos, se procedió a integrar
dicha información con la base de datos play-by-play. Para ello, los
nombres de los jugadores fueron previamente normalizados con el objetivo
de reducir posibles inconsistencias de formato (espacios adicionales,
diferencias menores de escritura), y luego se utilizó la correspondencia
jugador--equipo para asignar a cada acción del \emph{play-by-play} el
equipo del jugador involucrado.

De este modo, cada acción queda asociada explícitamente a uno de los dos
equipos del partido, aun cuando dicha información no se encuentre
disponible de forma directa en el registro original del
\emph{play-by-play}. Este procedimiento permite identificar que
jugadores estarán en cancha para cada uno de los equipos, permitiendo
construir los quintetos en cada momento del partido.

\paragraph{Identificación de los quintetos en cancha al momento de la
acción}\label{identificaciuxf3n-de-los-quintetos-en-cancha-al-momento-de-la-acciuxf3n}

A partir de la nueva tabla generada, se procedió a construir la
presencia de los jugadores en cancha a lo largo del partido. Para ello,
se identificaron las acciones asociadas a sustituciones, distinguiendo
entre entradas (acción = ENTRA A PISTA) y salidas de jugadores (acción =
ABANDONA LA PISTA).

\begin{figure}

{\centering \includegraphics[width=0.4\linewidth]{img/imagen sustitucion} 

}

\caption{Ejemplo de secuencia de acciones con sustitución}\label{fig:unnamed-chunk-4}
\end{figure}

\newpage

El conjunto de jugadores en cancha se actualiza a medida que se recorren
las acciones del partido, reiniciándose al comienzo de cada cuarto. De
este modo, para cada instante del juego se obtiene el conjunto de
jugadores activos en cancha.

Utilizando la tabla jugadores-equipos construida previamente, cada
jugador en cancha fue asignado al equipo local o visitante, permitiendo
reconstruir los quintetos de ambos equipos en cada acción del partido.
Este paso resulta esencial para caracterizar el contexto en el que
ocurre cada acción, ya que al conocer el equipo al que pertenece el
jugador que realiza la acción podemos inferir simultáneamente los
jugadores ofensivos y defensivos presentes en cancha en cada momento.

Una vez reconstruidos los quintetos en cancha para cada acción, se
incorporó un control de consistencia basado en la cantidad total de
jugadores activos. Dado que una situación de juego regular involucra
diez jugadores (cinco por equipo), se podrán identificar aquellas
acciones en las que la cantidad total de jugadores en cancha fuera
superior o inferior a dicho umbral.

\paragraph{Detección de errores en la carga de los
datos}\label{detecciuxf3n-de-errores-en-la-carga-de-los-datos}

Un factor importante a tener en cuenta es que los datos que se
encuentran en el \emph{play-by-play} son cargados de manera manual por
las mesas de control a medida que se desarrolla el partido. Esta tarea
se realiza en el contexto altamente dinámico de un partido, donde las
acciones se suceden con rapidez, lo que puede dar lugar a ciertos
errores en la registración de eventos. En consecuencia, el conjunto de
datos original puede contener inconsistencias que deben ser
identificadas y corregidas antes de su utilización en este trabajo.

Para intentar detectar posibles errores en la conformación de los
quintetos en cancha, se construyó , tal como se menciono anteriormente,
una variable que contabiliza la cantidad de jugadores presentes en cada
quinteto en cada una de las acciones del partido. Dicha variable debería
contabilizar sistemáticamente el valor cinco (5) para ambos equipos en
todas las acciones registradas.

De esta manera se detectó un caso puntual en el partido de ZARATE BASKET
vs ATENAS (C), donde el quinteto local registraba 6 jugadores en cancha.
Esta inconsistencia se observó desde el inicio del segundo cuarto (con
10:00 minutos restantes) hasta el minuto 4:42 del mismo período.

Para subsanar este inconveniente, se realizó una revisión manual del
play-by-play original disponible en la página web y se pudo identificar
que al comienzo del segundo cuarto se había registrado incorrectamente
el ingreso de seis jugadores en lugar de cinco, lo cual constituye un
error atribuible a la mesa de control. Una vez identificada esta
problemática, se reviso el contexto de las siguientes jugadas y pudo
determinarse que el jugador MERCHANT, EDGAR HENRY no se encontraba
realmente en cancha durante ese tramo del partido.

En consecuencia, se procedió a corregir el error eliminando al jugador
mencionado de la variable correspondiente al quinteto local en las
acciones afectadas, así como ajustando el conteo de jugadores en cancha.
De esta manera, se restauró la consistencia lógica de los quintetos y se
garantizó la validez de la información utilizada en las etapas
posteriores del análisis.

\section{Metodología}\label{metodologuxeda}

\section{Resultados}\label{resultados}

\section{Conclusión}\label{conclusiuxf3n}

\section{Discusión}\label{discusiuxf3n}

\newpage

\section{Bibliografía}\label{bibliografuxeda}

Agresti, A. (2013). \emph{Categorical Data Analysis} (3rd ed.). Wiley
Series in Probability and Statistics. Wiley, New York.

Damoulaki, A., Ntzoufras, I. \& Pelechrinis, K. (2025). \emph{Lasso
multinomial performance indicators for in-play basketball data.} Comput
Stat 40, 2157--2181.

Dunn, P. K., \& Smyth, G. K. (2018). \emph{Generalized linear models
with examples in R. Springer.}

Hvattum, L. M. (2019). \emph{A comprehensive review of plus--minus
ratings for evaluating individual players in team sports.} International
Journal of Computer Science in Sport, 18, 1--23.

Ilardi, S. (2007). \emph{Adjusted Plus-Minus: An idea whose time has
come.} Blog: 82games. \url{http://www.82games.com/ilardi1.htm}

Ilardi, S., \& Barzilai, A. (2008). \emph{Adjusted Plus-Minus ratings:
New and improved for 2007--2008.} Blog: 82games.
\url{http://www.82games.com/ilardi2.htm}

James, G., Witten, D., Hastie, T., \& Tibshirani, R. (2021). \emph{An
Introduction to Statistical Learning with Applications in R} (ISLR2).

Kubatko, J., Oliver, D., Pelton, K., \& Rosenbaum, D. (2007). \emph{A
starting point for analyzing basketball statistics.} Journal of
Quantitative Analysis in Sports, 3, Article 1.

Rosenbaum, D. (2004). \emph{Measuring how NBA players help their teams
win.} Blog: 82games. \url{http://www.82games.com/comm30.htm}

Sæbø, O., \& Hvattum, L. (2019). \emph{Modelling the financial
contribution of soccer players to their clubs.} Journal of Sports
Analytics, 5, 23--34.

Sill, J. (2010). \emph{Improved NBA adjusted +/-- using regularization
and out-of-sample testing.} Proceedings of the 2010 MIT Sloan Sports
Analytics Conference.

Tibshirani, R. (1996). \emph{Regression shrinkage and selection via the
lasso.} Journal of the Royal Statistical Society: Series B
(Methodological), 58(1), 267--288.

Zou, H., \& Hastie, T. (2005). \emph{Regularization and Variable
Selection via the Elastic Net.} Journal of the Royal Statistical
Society: Series B, 67(2), 301--320.

\end{document}
